% ===========================================================================
\chapter{Not doing the same work twice}\label{ch:seminaive}

The naive evaluator is correct, and wasteful. Each round applies every rule
to every fact, so each round rediscovers everything every earlier round
found. Take a chain of towns, each linked to the next, and the rules for
routes along it:

\begin{shell}
route(X, Y) :- link(X, Y).
route(X, Z) :- route(X, Y), link(Y, Z).
\end{shell}

In the first round the routes of one link appear; in the second, the routes
of two links; and so on, one link longer each round. But every round, the
second rule joins \emph{all} the routes found so far with the links, and
all but the newest of those joins only find routes already known.

\section{The idea}

A fact that is new in this round needs a reason that is new too. If every
fact a rule's body used was already known last round, the rule already
derived its head last round. So, after the first round, each rule only
needs to be joined in the ways that use \emph{at least one fact the
previous round added}.

That is easy to arrange. Keep the facts the previous round added in a small
database of their own, the \emph{delta}. For a rule with several positive
atoms, join it once for each of them, reading that atom's facts from the
delta and the rest from the whole database. Every derivation that uses a
new fact is found by at least one of those joins. This is called
\emph{semi-naive} evaluation.

\section{The semi-naive evaluator}

The join already asks \code{source} where an atom's facts come from. The
semi-naive evaluator overrides it: the atom at position \code{delta_at}
reads from the delta.

\begin{append}{evaluate.nex}
-- Is `l` a positive literal, one that facts can be found for?
function is_positive(l: Literal): Boolean
do
  match l of
    Positive then result := true
    _        then result := false
  end
end
\end{append}

\begin{continued}{evaluate.nex}
-- After the first round, join each rule only against the facts that the
-- previous round added: a new fact needs at least one new premise.
class Semi_Naive_Evaluator inherit Evaluator
create
  make
  do
    delta := create Database.make
    delta_at := -1
  end
feature
  -- The facts the previous round added.
  delta: Database
  -- Which literal of the body reads `delta`; -1 for none.
  delta_at: Integer

  source(i: Integer, atom: Atom, b: Bindings, db: Database): Array[Atom]
  do
    if i = delta_at then
      result := delta.candidates(atom, b)
    else
      result := db.candidates(atom, b)
    end
  end

  run(rules: Array[Rule], db: Database)
  do
    -- The first round has nothing to compare with: every rule, every fact.
    count_round
    delta_at := -1
    let fresh := create Database.make
    across rules as r do
      apply(r, db, fresh)
    end
    from until fresh.count = 0 do
      count_round
      delta := fresh
      fresh := create Database.make
      across rules as r do
        from let i := 0 until i = r.body.length do
          if is_positive(r.body.get(i)) then
            delta_at := i
            apply(r, db, fresh)
          end
          i := i + 1
        end
      end
    end
    delta_at := -1
  end

  -- Derive what `r` gives now, adding the new facts to `db` and `fresh`.
  apply(r: Rule, db: Database, fresh: Database)
  do
    across solutions(r, db) as b do
      let f := substitute(r.head, b)
      if db.add_derived(f, r, b) then
        fresh.add(f)
      end
    end
  end
end
\end{continued}

Notice what the new class does not have: a postcondition. It inherits
\code{run}'s from \code{Evaluator}, so every run of the semi-naive evaluator
is checked, by the plain apply-every-rule test of \code{is_fixpoint}, to
have found everything the naive one would. The faster algorithm is held to
the simple one's standard, at no cost in code.

\section{Trying it out}

The chain in \code{chain.dl} has thirty towns. Both evaluators should find
the same 464 facts; the question is how much work each does.

\begin{tryit}{try_seminaive.nex}
intern Parser
intern Evaluate
intern io/Path

function evaluated(p: Datalog_Program, e: Evaluator): Database
do
  result := create Database.make
  across p.facts as f do
    result.add(f)
  end
  e.run(p.rules, result)
end

match parse_program(create Path.make(program.working_directory + "/chain.dl").read_text) of
  Ok as ok then
    do
      let naive := create Naive_Evaluator.make
      let semi := create Semi_Naive_Evaluator.make
      let a := evaluated(ok.value, naive)
      let b := evaluated(ok.value, semi)
      console.print_line("naive:       " + a.count.to_string + " facts, "
        + naive.rounds.to_string + " rounds, " + naive.attempts.to_string + " facts looked at")
      console.print_line("semi-naive:  " + b.count.to_string + " facts, "
        + semi.rounds.to_string + " rounds, " + semi.attempts.to_string + " facts looked at")
    end
  Err as e then console.print_line(e.error)
end
\end{tryit}
\tryoutput{try_seminaive}

The same facts, in the same number of rounds, for about a tenth of the work.
The saving grows with the length of the chain, because the naive evaluator
re-joins every route it knows in every round.
